\documentclass[a4paper,11pt]{article}
\usepackage[left=2cm,text={17cm,24cm},top=3cm]{geometry}
\usepackage[czech]{babel}
\usepackage[T1]{fontenc}
\usepackage[utf8]{inputenc}
\usepackage{times}


\begin{document}

    \begin{titlepage}
        \begin{center}
            \textsc{\Huge{Vysoké učení technické v Brně}} \\
            \vspace{0.5 em}
            \textsc{\huge{Fakulta informačních technologií}} \\
            \vspace{\stretch{0.382}}
            \LARGE{Typografie a publikování\,--\,4. projekt} \\
            \Huge{Bibliografické citace}
            \vspace{\stretch{0.618}}
        \end{center}
        \Large{\today \hfill Jan Černoch}
    \end{titlepage}

\clearpage
\section*{Mikrotypografie v digitální sazbě}

Mikrotypografie označuje soubor jemných typografických úprav, které zlepšují vizuální kvalitu textu. 
Tyto úpravy se zaměřují především na detaily, jako jsou mezery mezi znaky, zarovnání okrajů nebo použití ligatur. 
Na rozdíl od makrotypografie, která řeší strukturu dokumentu, mikrotypografie pracuje s jednotlivými znaky a jejich optickým vnímáním. 
Podle Wagnera mají právě tyto drobné úpravy zásadní vliv na čitelnost delších textů \cite{wagner}. 

Jedním z nejznámějších principů mikrotypografie je optické vyrovnání okrajů. 
Některé znaky, například uvozovky nebo pomlčky, mohou mírně přečnívat přes okraj sazby. 
Tím se docílí vizuálně rovného okraje, který působí přirozeněji než striktně zarovnaný blok textu. 
Tento princip je využíván zejména v profesionální knižní sazbě \cite{bringhurst}. 

Další důležitou technikou je jemné rozšiřování znaků, označované jako font expansion. 
Tato metoda umožňuje upravit šířku znaků o velmi malé hodnoty, čímž se zlepší zalomení odstavce. 
Rozdíl není pro čtenáře patrný, ale výrazně ovlivňuje kvalitu sazby. 
Digitální sázecí systémy tuto techniku implementují automaticky \cite{microtype}. 

Mikrotypografie úzce souvisí také s kerningem. 
Kerning upravuje mezery mezi konkrétními dvojicemi znaků, které by jinak působily opticky nevyváženě. 
Typickým příkladem jsou kombinace písmen „AV“ nebo „To“. 
Správně nastavený kerning zlepšuje plynulost textu a zabraňuje vzniku rušivých mezer \cite{root}. 

Dalším prvkem jsou ligatury. 
Ligatury spojují vybrané dvojice znaků do jednoho glyfu, čímž se zlepšuje vizuální kontinuita textu. 
Moderní formáty písem umožňují definovat ligatury pomocí OpenType funkcí \cite{opentype}. 
Použití ligatur je běžné zejména v kvalitní knižní sazbě. 

V digitálním prostředí se mikrotypografie prosadila především díky systému \LaTeX. 
Balíček microtype umožňuje automatickou aplikaci protrusion a font expansion. 
Tyto úpravy zlepšují zalamování řádků a celkový vzhled dokumentu \cite{latexmicro}. 
Optimalizační algoritmus pro zalamování řádků navržený Knuthem patří mezi nejkvalitnější dostupné přístupy \cite{knuth}. 

Problematika mikrotypografie je diskutována i v odborných časopisech. 
Periodikum TUGboat pravidelně publikuje články věnované kvalitě sazby \cite{tugboat}. 
Například Mittelbach popisuje rozšíření systému \LaTeX\ o mikrotypografické funkce \cite{mittelbach}. 
Další práce se zabývá praktickým využitím těchto technik v dokumentech \cite{hoenig}. 

Mikrotypografie je rovněž předmětem kvalifikačních prací. 
Například diplomová práce na VUT analyzuje implementaci mikrotypografie v systému \LaTeX\ a její vliv na kvalitu sazby \cite{vutthesis}. 
Další práce se zaměřuje na využití typografických pravidel v elektronických dokumentech \cite{cvutthesis}. 
Výsledky těchto studií potvrzují, že i malé úpravy mohou výrazně zlepšit čitelnost textu. 

Bringhurst uvádí, že „typografie existuje proto, aby uctila obsah“ \cite{bringhurst}. 
Mikrotypografie tento princip naplňuje na úrovni jednotlivých znaků. 
Správně aplikované mikrotypografické úpravy zvyšují estetickou kvalitu dokumentu i komfort čtení. 
V digitální sazbě proto představují důležitou součást profesionální typografie. 

\clearpage
\bibliographystyle{czechiso}
\bibliography{proj4}

\end{document}
